% Source: round6/symmetry.tex
\section{Coordinate permutation symmetry prevents entropy attainment}
\label{s12:sec:coordinate-symmetry}
Permutation symmetry among the \emph{bit coordinates} is different
from exchangeability among tensor copies in the earlier quantum
formulation. Here we prove that imposing full coordinate symmetry
destroys the desired entropy concentration once there are at least
five bits. The proofs are elementary and do not use a de Finetti approximation.

Let $w$ be an unbiased pairwise-independent law on $d\ge2$ sign bits,
invariant under all coordinate permutations. Write $K$ for the number
of positive signs and $a_k=\Pr[K=k]$. Then each atom in the $k$th
slice has probability $a_k/\binom dk$, and
\[
 \mathbb E K=d/2,\quad \operatorname{Var}K=d/4,\quad
 C(w)=\sum_{k=0}^d\frac{a_k^2}{\binom dk}.
\]

\begin{theorem}[Collision cost of full coordinate symmetry]
\label{s12:thm:symmetry-collision}
For $d\ge7$,
\[
 C(w)\le B_d:=\frac1{d^2}+\frac d{(d-2)^4}
                    +\frac2{d(d-1)}<\frac4{d^2}.
\]
The maximal collision probability among such symmetric laws is
$\Theta(d^{-2})$. For every $d\ge5$, no such law attains the
residue-sensitive optimum
$1/N$, $N=4\lceil(d+1)/4\rceil$.
For each $d\in\{2,3,4\}$ a symmetric law attaining that bound exists.
\end{theorem}
\begin{proof}
Put $S=2K-d$, so $\mathbb E S^2=d$. The total masses in the two
endpoint slices and their two neighbors satisfy
\[
 a_0+a_d\le1/d,\qquad a_1+a_{d-1}\le d/(d-2)^2.
\]
Their collision contributions are at most $1/d^2$ and
$d/(d-2)^4$. All other slices have multiplicity at least
$\binom d2$, so their total collision contribution is at most
$1/\binom d2$. This proves the first bound. For $d\ge7$,
$2d/(d-1)\le7/3$, while $d^3/(d-2)^4$ decreases with $d>2$
and its value at seven is $343/625<2/3$. Thus $d^2B_d<4$.
Since $d^2>4(d+4)$ for $d\ge7$ and $N\le d+4$, we have
$C(w)<1/(d+4)\le1/N$.

For $d=5$, equality in the residue theorem would make $w$ uniform
on eight atoms. A permutation-invariant support is a union of slices,
whose sizes are $1,5,10,10,5,1$. No subcollection sums to eight.
For $d=6$, the sizes are $1,6,15,20,15,6,1$. An eight-atom union must
consist of both endpoint slices and one neighboring size-six slice.
Its uniform law has nonzero mean $\mathbb E S=\pm3$, so cannot be
unbiased. For $d=2$ take the uniform cube. For $d=3,4$ take the uniform
even-parity cube; each pair of coordinates is uniform and the support
sizes are respectively four and eight.

For the matching order-of-growth lower bound when $d$ is odd, put
mass $1/(d+1)$ on the positive vector and distribute the remaining
mass uniformly on $S=-1$. This is unbiased and pairwise independent
by its first two $S$ moments, and has collision probability at least
$1/(d+1)^2$. When $d$ is even, put mass $1/(2d)$ on each endpoint
and distribute the remaining mass uniformly on $S=0$. Its first two
$S$ moments again give the required independence, and its collision
probability is at least $1/(2d^2)$. Both constructions are symmetric.
\end{proof}

\paragraph{Formalization of the collision theorem.}
Seven independent \texttt{CollisionChallenge} targets cover the full
quantitative envelope, its strict inverse-square bound, an actual
maximizer and its formal $\Theta(d^{-2})$ rate, all nonattainment cases,
and exact attainment at dimensions two, three, and four.
The formal proof uses exact coordinate orbits and Hamming moments,
explicit odd and even slice mixtures, finite slice counts for the two
exceptional nonattainment cases, and compactness for maximum attainment.
No definition placeholders occur in the independent challenge.

\begin{theorem}[Shannon entropy is almost maximal]
\label{s12:thm:symmetry-shannon}
Every such symmetric pairwise-independent law satisfies
\[
 H(w)\ge(d-1)\log2-\log(d+1).
\]
Thus full coordinate symmetry forces Shannon entropy to within
$O(\log d)$ of the maximum $d\log2$, rather than the desired
order $\log d$.
\end{theorem}
\begin{proof}
Conditional uniformity on slices gives
$H(w)=H(K)+\mathbb E\log\binom dK$.
Let $h(p)=-p\log p-(1-p)\log(1-p)$. For $p=k/d$, the binomial
mass at $k$ is a mode, so it is at least $1/(d+1)$. Hence
\[
 \log\binom dk\ge d h(k/d)-\log(d+1).
\]
The elementary bound $h(p)\ge4(\log2)p(1-p)$ follows, for example,
from the convergent expansion
\[
 h(p)=\log2-\sum_{j\ge1}
           \frac{(2p-1)^{2j}}{2j(2j-1)}:
\]
the positive coefficients sum to $\log2$, and every even power is
at most $(2p-1)^2$. Continuity handles the endpoints.
Finally $\mathbb E[K(d-K)]=d(d-1)/4$, giving
$\mathbb E[d h(K/d)]\ge(d-1)\log2$; discard $H(K)\ge0$.
\end{proof}

\paragraph{Formalization of the Shannon bound.}
The independent \texttt{SymmetricChallenge} states the full theorem for
an actual nonnegative mass-one law and its coordinate permutation
invariance and first two sign moments. The Lean proof classifies
coordinate orbits, counts their binomial slice sizes, and proves the
binomial-mode estimate. It proves the quadratic binary-entropy bound
by integrating a pointwise harmonic inequality, providing an alternative
to the series argument above. All endpoint cases are included.

This does not rule out symmetric designs under smaller groups, nor
prove nonexistence of Hadamard matrices. It rules out a specific
attainment strategy: averaging over \emph{all} bit permutations while
hoping to retain minimal collision or Shannon entropy. The sought
low-entropy construction must break that symmetry for $d\ge5$.

